%\chapter{Discusión de resultados}


%El enfoque seguido se basó primero en examinar
%las matrices de confucion y de estas medidas fue posible inferir que el modelo que mejor desepeño tiene en ambos sets de datos en el modelo You Only Look Once (YOLO) \cite{Redmon2016} el cual se usara como piedra angular para la construcción de la arquitectura federada.


%Con el modelo base elegido se construyo una arquitectura federada sencilla en donde se simularon dos clientes y el servidor en donde la agregación de los pesos sera efectuada. Los resultados presentados en la seccion anterior son un conjunto de metricas calculadas en el proceso de validacion y testing del modelo bajo los mencionados algoritmos de agregacion para los modelos federados. El primer resultado directo obtenido a partir de una inspeccion de la forma de las curvas por clase y su comportamiento global, identificando puntos de inflexión o máximos que sugieran umbrales de confianza óptimos. Es que no existe evidencia de una diferencia estadistica significativa entre los algoritmos de agregacion en este set de datos. 

%Posteriormente, se compararon los resultados por clase con la curva promedio (representada en azul en cada gráfico) enfocandonos ya en el comportamiento concreto de las curvas  para tener una idea clara del rendimiento global del modelo y las desviaciones por clase. Finalmente, se interpretaron los valores del mAP@0.5 y otros indicadores directos de desempeño, como la precisión o recall máximo alcanzado, relacionándolos con el tipo de clase y posibles dificultades inherentes a su detección. A partir de esto se 

%\section{Observaciones por gráfico}

%En primer lugar, la curva F1-Confidence muestra cómo varía el balance entre precisión y recall en función del umbral de confianza. Se observa que el modelo logra un máximo F1-score promedio de aproximadamente 0.85 cuando el umbral de confianza es 0.275. Este valor representa el mejor punto de equilibrio global, donde ni la precisión ni el recall predominan, sino que ambos contribuyen equitativamente al rendimiento general. Clases como \textit{Mascara\_Codigos} y \textit{Mascara\_ColorChecker} muestran curvas notablemente altas y planas, lo cual indica un comportamiento robusto a lo largo de distintos niveles de confianza. En contraste, la clase \textit{Mascara\_Encabezados} presenta un F1-score bajo y decreciente, lo cual sugiere que el modelo tiene dificultades tanto para identificar correctamente estas instancias como para evitar falsos positivos.

%Al observar la curva Recall-Confidence, se confirma el comportamiento esperado de que el recall disminuye progresivamente con umbrales más altos de confianza. En el punto extremo de umbral cero, el modelo logra un recall promedio cercano al 0.99, lo que significa que es capaz de detectar casi todas las instancias verdaderas si se permiten predicciones de baja certeza. Nuevamente, las clases de mejor desempeño mantienen altos valores de recall en rangos amplios de confianza, mientras que la clase \textit{Encabezados} muestra un descenso abrupto desde niveles muy bajos de confianza, lo que delata una gran cantidad de falsos negativos incluso con umbrales permisivos.

%La tercera gráfica, Precision-Confidence, revela cómo mejora la precisión a medida que se aumenta el umbral de confianza, sacrificando recall como ya fue visto. El modelo alcanza una precisión promedio de 1.0 cuando la confianza es 1.0, lo cual es natural ya que se están considerando únicamente las predicciones más seguras. Lo más relevante aquí es que el resto de las clases, excepto \textit{Encabezados}, se mantienen en rangos altos de precisión incluso con umbrales bajos, lo que evidencia un modelo confiable en la mayoría de los casos. Sin embargo, la curva de la clase \textit{Mascara\_Encabezados} vuelve a destacar negativamente, ya que parte de una precisión muy baja, por debajo de 0.4, lo que sugiere una alta tasa de falsos positivos.

%Finalmente, la curva de Precisión vs Recall es probablemente la más ilustrativa para visualizar el comportamiento general del modelo. La métrica mAP@0.5 promedio es de 0.886, lo cual representa un rendimiento muy bueno en términos de detección. Las clases \textit{ColorChecker} y \textit{Codigos} tienen un desempeño casi perfecto, con valores superiores a 0.99. Por otro lado, la clase \textit{Encabezados} vuelve a presentar el rendimiento más bajo, con un mAP@0.5 de apenas 0.509, muy por debajo del resto. Esto refuerza la idea de que esta clase específica representa el mayor desafío del modelo y es la principal responsable de la brecha entre el rendimiento ideal y el real.
\chapter{Discusión de resultados}

\section{Comparación con la literatura}

Los resultados obtenidos confirman que la arquitectura \textit{You Only Look Once} (YOLO) \citep{Redmon2016} superó en desempeño a las alternativas evaluadas (RCNN y SSD), alcanzando un mAP@0.5 promedio de 0.886. Este hallazgo concuerda con lo reportado en la literatura, donde YOLO ha mostrado ventajas en detección en tiempo real gracias a su reformulación del problema como regresión directa desde píxeles a coordenadas y categorías, eliminando la necesidad de etapas de propuestas de regiones \citep{Ren2015, Simonyan2014}. Asimismo, la robustez del modelo en clases visualmente regulares como \textit{ColorChecker} y \textit{Códigos} es coherente con estudios que demuestran que la uniformidad estructural de los objetos facilita la generalización \citep{Goodfellow2016, powers2020evaluation}.

En cuanto al aprendizaje federado, la implementación de FedAvg permitió mantener un alto desempeño sin comprometer la privacidad de los datos, tal como lo proponen \citet{McMahan2017} y \citet{bonawitz2016practicalsecureaggregationfederated}. El hecho de que no se observaran diferencias estadísticamente significativas entre algoritmos de agregación (FedAvg, FedProx, FedOpt, FedSVRG) coincide con estudios que indican que, en escenarios de heterogeneidad moderada, FedAvg es suficiente para garantizar estabilidad y convergencia \citep{mansour2022federated, ghosh2020efficient}. No obstante, la literatura también advierte que en contextos con alta variabilidad de los datos, algoritmos más robustos pueden ofrecer ventajas \citep{blanchard2017machine, yin2018byzantine}.

\section{Limitaciones del estudio}

A pesar de los buenos resultados globales, el modelo presenta debilidades claras en la detección de la clase \textit{Encabezados}, cuyo mAP@0.5 apenas alcanzó 0.509. Este bajo desempeño puede atribuirse a factores como: (i) la escasa representación de esta clase en el conjunto de entrenamiento, (ii) la alta variabilidad en tipografías, tamaños y posiciones, y (iii) posibles inconsistencias en las anotaciones. Estas limitaciones se alinean con lo señalado por \citet{sokolova2009systematic} y \citet{huang2007loss}, quienes destacan la dificultad de modelos de clasificación en el tratamiento de clases poco frecuentes o de elevada variabilidad intra-clase.

Adicionalmente, aunque el aprendizaje federado mostró un buen desempeño en este caso, su implementación enfrenta retos prácticos en escenarios reales. Entre ellos destacan la heterogeneidad de hardware en los clientes, la necesidad de sincronización eficiente y los riesgos de convergencia lenta cuando se trabaja con bases de datos mucho más desbalanceadas \citep{reddi2020adaptive}. Por lo tanto, los resultados deben interpretarse como una validación inicial más que como una garantía absoluta de rendimiento en condiciones operativas de mayor complejidad.

\section{Implicaciones para herbarios digitales y perspectivas futuras}

El hallazgo de un umbral de confianza óptimo cercano a 0.275, en el cual se maximiza el F1-score promedio (0.85), resulta particularmente relevante para aplicaciones prácticas de curaduría digital. Este balance entre \textit{precision} y \textit{recall} es crucial en contextos donde los falsos negativos implican pérdida de información taxonómica valiosa. Estudios como los de \cite{chen2020fedbe} resaltan la necesidad de calibrar dinámicamente los umbrales según las prioridades de cada caso de uso, lo cual también sería aplicable a herbarios digitales de diferentes tamaños y recursos.

En términos prácticos, los resultados respaldan el uso de YOLO combinado con aprendizaje federado como una herramienta escalable para la automatización de la curaduría de herbarios. Este enfoque permitiría la integración de colecciones distribuidas sin necesidad de centralizar los datos, favoreciendo la colaboración interinstitucional y garantizando la privacidad de la información, como ya se ha explorado en otros dominios sensibles como la salud y las finanzas \cite{bonawitz2016practicalsecureaggregationfederated,McMahan2017}. 

Finalmente, el estudio abre líneas de investigación futuras: (i) incrementar el tamaño y balance de las clases mediante \textit{data augmentation}, (ii) mejorar la calidad de anotaciones con procesos de curación manual, (iii) explorar arquitecturas modernas con mecanismos de atención espacial \citep{He2016}, y (iv) probar algoritmos de agregación robusta que permitan enfrentar escenarios de datos altamente heterogéneos \citep{mansour2022federated,pillutla2019robust}. De este modo, la propuesta aquí desarrollada puede escalarse y consolidarse como una solución integral para la gestión digital de colecciones botánicas a gran escala.